\section{Introdução}
\label{introducao}

A evasão é um dos problemas mais persistentes do ensino superior: cada estudante que abandona o curso representa recurso público não convertido em formação e, sobretudo, uma trajetória pessoal interrompida. Com a ampliação do acesso e a disponibilidade crescente de dados acadêmicos, a Mineração de Dados Educacionais passou a apoiar a identificação precoce de estudantes em risco \citep{romero2020}. A literatura clássica, porém, já descreve a evasão como um processo longitudinal e multifatorial, resultante da interação entre atributos de entrada, integração acadêmica e condições externas \citep{tinto1975}, e não como efeito de uma causa única.

Este trabalho delimita-se a uma instituição pública de ensino superior portuguesa, cujos registros de 4.424 estudantes de 17 cursos de graduação, ingressantes entre 2008/09 e 2018/19, foram disponibilizados por \citet{realinho2022}. A unidade de análise é o estudante, descrito por atributos demográficos, socioeconômicos, de ingresso e de desempenho nos dois primeiros semestres, e o desfecho observado é \textit{evasão}, \textit{matrícula ativa} ou \textit{graduação}.

\textbf{Problema e lacuna.} Instrumentos tradicionais de monitoramento analisam variáveis isoladamente (média, renda, distância, reprovações) e respondem ``quanto'' cada fator pesa, mas não ``como'' os fatores se combinam. Dois estudantes com o mesmo valor de um indicador podem seguir caminhos distintos. Classificadores supervisionados atingem boa acurácia neste conjunto de dados \citep{realinho2022}, porém produzem uma pontuação por indivíduo sem tornar explícita a estrutura de perfis semelhantes e de co-ocorrência de fatores. Faltam, no contexto educacional, modelos relacionais que tornem essa estrutura observável e interpretável pela gestão acadêmica, lacuna também apontada por \citet{porto2025} para perfis socioeconômicos.

\textbf{Questão de pesquisa.} \textit{Como a representação dos estudantes e de seus fatores como grafos revela perfis e combinações de fatores associados à evasão e ao sucesso acadêmico, para além das análises isoladas de variáveis?} Dela decorrem as hipóteses:
\textbf{H1}~--~em um grafo de similaridade construído \textit{sem} o desfecho, estudantes semelhantes tendem a compartilhar o mesmo desfecho (homofilia);
\textbf{H2}~--~a detecção de comunidades isola perfis com taxas de evasão muito distintas da média institucional (32{,}1\%);
\textbf{H3}~--~um escore de risco derivado da vizinhança no grafo supera o melhor indicador isolado e é competitivo com um modelo tabular multivariado; e
\textbf{H4}~--~na rede de co-ocorrência de fatores, os fatores associados à evasão formam um agrupamento coeso e distinto dos associados ao sucesso.

\textbf{Justificativa.} Do ponto de vista científico, o trabalho propõe e compara dois modelos de grafo complementares (estudante--estudante e fator--fator) para dados educacionais tabulares, seguindo o processo CRISP-NET \citep{acuna2025}. Do ponto de vista institucional, coordenações e pró-reitorias precisam decidir, com recursos limitados, \textit{quem} acompanhar e \textit{quais combinações} de condições merecem ações de permanência; a estrutura relacional oferece grupos acionáveis em vez de listas de variáveis.

\textbf{Objetivos.} O objetivo geral é modelar e analisar os registros acadêmicos como redes para identificar perfis de risco e padrões associados aos desfechos. Os objetivos teóricos são: (\textit{i}) formalizar um grafo de similaridade multirrelacional e uma rede de fatores; e (\textit{ii}) avaliar quantitativamente sua estrutura por modularidade, assortatividade e validação cruzada. Os objetivos práticos são: (\textit{iii}) implementar um protótipo reprodutível de análise e exploração; e (\textit{iv}) caracterizar perfis e combinações de fatores que apoiem ações de acompanhamento e intervenção.

O documento está organizado da seguinte forma: a Seção~\ref{referencial} apresenta os conceitos e trabalhos relacionados; a Seção~\ref{metodologia} descreve o dataset, os modelos e o método de solução; a Seção~\ref{resultados} apresenta e discute os resultados; e a Seção~\ref{conclusoes} conclui o trabalho.
